% show abstract on the title page
\begin{abstract}

Population ageing has raised concerns about the future of aggregate labour supply, but its effects depend not only on the rising share of older individuals, but also on immigration, participation, and hours worked. This article estimates aggregate labour supply in Spain using Labour Force Survey microdata and assesses the contribution of each component to changes between 2005 and 2025, using a continuous change decomposition framework. Over this period, aggregate labour supply increased by 2.18 million FTE workers. Population ageing among the Spanish-born population reduced labour supply by 2.55 million, a decline that was more than compensated for by immigration, which added 3.04 million FTE workers. However, higher labour-market participation among prime-age and older Spanish-born individuals contributed a similar amount as immigration (+2.82 million), even though declining hours worked mitigated the gain (-1.60 million).

\vspace{0.7em}\par
\textbf{Keywords}: Labour Supply, Immigration, Population Ageing, Decomposition Models.
\end{abstract}

